% Tapis de Quillen, with Quillen's offprint and preprint (cote 162-5, typed part): archive cote 162-5, site page 17 (PDF page 18).
% typescript. Draft compared: none.
% AI-transcribed by Claude Opus 5.5 (claude-opus-5-5, adaptive thinking, effort high), Message Batches API, batch t162_full (msgbatch_01UAh3zyWPd454AeDkhnmQ33), 2026-11-13; not yet checked by eye.
\scanpage{162-5/17}
\mspage{16}

Le théorème de Thom est alors qu'un $x \in H^{\cdot}(X)$ est dans $S(X)$ sss il est
\xout{\illegible{1}}\inserted[typed]{fixe} par G, i.e. satisfait $A^{+}.\ x = 0$. (Pourquoi cette condition est-el%
% NOTE: the right edge of the scan cuts the ends of several typed lines
le déjà nécessaire ?) Une autre façon de le voir est de dire que le
foncteur $H^{\cdot}$ considéré comme allant de la catégorie $\underline{S}$ précédente dans
la catégorie des \xout{\unsure{vectoriels}} représentations de G dans des vectoriels
sur $\underline{F}_2$ , est pleinement fidèle. \xout{Comme G est} Quand on prend l'enveloppe
pseudo-abélienne, on trouve une sous-catégorie pleine de la catégorie des
représentations de G, mais il est douteux qu'on trouve une catégorie abé-%
lienne, voire \unsure{la|ta} catégorie de toutes les représentations de G de dimensi
finie sur $\underline{F}_2$ . (Quelle est la structure de G, qui est prolisse; d'après
\margin{\illegible{20}}
Quillen, il serait pro-\unsure{nilpotent} ....). Ici G bien sûr joue le rôle d'u
groupe de Galois motivique, mais contrairement à la situation familière,
il semble qu'il soit fortement non réductif (et d'ailleurs le corps de
\hand{$\triangleright$} \unsure{définition} dans tout ceci est $\underline{F}_2$ , et non $\underline{Q}$ ! ).

Quillen\deleted{\unsure{t}} donne une généralisation de la construction d'une caté-%
gorie universelle $\underline{B}$ au cas où on travaille avec des variétés pas nécessa
rement propres (mais $f_{.}$ n'étant supposé défini que pour f propre), et
quand on introduit une condition d'orientation \inserted[typed]{stable} sur les \uline{morphismes}
envisagés , via un morphisme $H \to BO(\infty)$, où H est un H-espace (commuta
tif, associatif et unitaire ?): une orientation de $f:X \to Y$ sera un
relèvement homotopique de $X \to BO$, défini par \xout{\illegible{1}} $T_X - f^{\cdot}(T_Y)$ , en $X \to H$
les éléments de $H(X,Y)$ seront définis par des diagrammes de morphismes
orientés $Z \to X$, $Z \to Y$, le deuxième étant propre, mudulo bords de tels.
En tant qu'ensemble, il se définit aussi en termes de la vari\hand{é}té XxY
munie simplement de la famille des supports formée des ensembles fermés
propres sur Y ..... \hand{Les $B^{H}_{m}(X)$ (X \unsure{compact disons}) peuvent}
\hand{se définir homotopiquement à la Thom, comme les groupes d'homotopie}
\hand{\unsure{stables d'espaces} \uline{Hom}$(X, MH(N-m))$ -----.}

\hand{cf exposé. Karoubi à \uline{Bourbaki}}

\margin{17}
